\chapter{LSD (overview)}
 
\section*{Definition and Overview}
Lysergic acid diethylamide (LSD), commonly known as "acid," is a semi-synthetic, highly potent psychedelic substance derived from lysergic acid, an alkaloid found in the ergot fungus (\textit{Claviceps purpurea}). Synthesized originally in 1938, LSD is classified as a classic hallucinogen. It acts primarily as an agonist at serotonin 5-HT2A receptors in the brain, inducing profound alterations in sensory perception, mood, thought processes, and self-awareness, collectively referred to as a "trip." LSD is a Schedule 9 controlled substance in Australia, indicating it has no established therapeutic use and is subject to strict prohibition, although modern research is investigating its potential under clinical supervision for treating conditions like treatment-resistant depression and anxiety.
 
\section*{Epidemiology}
According to national drug strategy household surveys, the lifetime prevalence of hallucinogen use (including LSD and psilocybin) among Australians aged 14 and over is approximately 10--12\%, with annual use generally low at around 1--2\%. LSD is most commonly used by young adults aged 18 to 29 years. Acute toxicity presentations to emergency departments are relatively uncommon compared to other illicit drugs like alcohol, stimulants, or opioids. When presentations occur, they are typically related to severe psychological distress ("bad trips") or accidental injury due to impaired judgment rather than direct physical organ toxicity.
 
\section*{Aetiology and Pathophysiology}
The physiological and psychological effects of LSD are mediated through specific receptor interactions in the central nervous system:
 
\begin{itemize}
    \item \textbf{Serotonin receptor agonism:} LSD is a partial agonist at serotonin 5-HT2A receptors, particularly in the prefrontal cortex, which plays a major role in mood, sensory processing, and cognition.
    \item \textbf{Dopamine and glutamate interaction:} it also interacts with dopamine D2 receptors and enhances glutamatergic transmission, contributing to its complex psychological effects.
    \item \textbf{Altered sensory gating:} LSD disrupts the thalamocortical gating system, allowing sensory information to bypass normal filters, resulting in synaesthesia (e.g., "hearing colors" or "seeing sounds") and vivid hallucinations.
    \item \textbf{Default mode network (DMN):} it alters connectivity in the default mode network, which is associated with ego dissolution and feelings of connectedness.
\end{itemize}
 
\section*{Clinical Features}
The effects of LSD typically begin within 30 to 90 minutes of ingestion and can last for 8 to 12 hours:
 
\begin{itemize}
    \item Hallucinations (mostly visual, such as moving geometric patterns, trailing images, or distorted shapes)
    \item Synaesthesia, altered passage of time, and heightened sensory awareness
    \item Psychological changes: intense mood swings, depersonalization, derealization, and ego dissolution
    \item Sympathomimetic features: mild tachycardia, pupillary dilation (mydriasis), mild hypertension, sweating, and tremors
    \item "Bad trip" features: acute panic, severe anxiety, paranoia, fear of losing control, and dysphoria
    \item Chronic or delayed effects: flashbacks (brief recurrences of drug-like experiences long after the drug has cleared, known as Hallucinogen Persisting Perception Disorder, or HPPD)
\end{itemize}
 
\section*{Differential Diagnosis}
\begin{itemize}
    \item \textbf{Acute stimulant toxicity (amphetamines, cocaine):} presents with sympathomimetic features and psychosis, but is characterized by prominent agitation, severe hypertension, and paranoia, with less emphasis on vivid visual hallucinations.
    \item \textbf{Primary psychotic disorders (schizophrenia, bipolar mania):} presenting with hallucinations and delusions, but typically lack the rapid onset, visual focus of hallucinations, and physical signs (like mydriasis) of LSD exposure.
    \item \textbf{Anticholinergic syndrome:} presents with visual hallucinations, mydriasis, and tachycardia, but features dry skin, dry mouth, urinary retention, and absent bowel sounds.
    \item \textbf{Encephalitis or delirium:} presenting with altered mental state, distinguished by fever, neck stiffness, or localising neurological signs.
\end{itemize}
 
\section*{When to Seek Medical Care}
Seek medical evaluation if an individual is experiencing severe agitation, panic, paranoia, or is behaving in a way that poses a danger to themselves or others.
 
\textbf{Call 000 (or local emergency services) for:}
\begin{itemize}
    \item Extreme, uncontrollable agitation, violence, or severe paranoia
    \item Seizures or loss of consciousness
    \item High fever (hyperthermia) combined with rapid heart rate
    \item Thoughts of self-harm or suicide during or after a trip
    \item Accidental trauma or injury sustained while under the influence
\end{itemize}
 
\section*{Diagnosis and Investigations}
\begin{itemize}
    \item \textbf{Clinical evaluation:} diagnosis is primarily based on history and typical clinical findings, as standard rapid urine drug screens do not detect LSD.
    \item \textbf{Specific toxicology tests:} liquid chromatography-mass spectrometry (LC-MS) can detect LSD in blood or urine, but results are rarely available in emergency settings.
    \item \textbf{ECG:} performed if tachycardia or chest pain is present.
    \item \textbf{Blood tests:} electrolytes, kidney function, and creatine kinase (CK) if severe agitation or prolonged restraint raises concern for rhabdomyolysis.
\end{itemize}
 
\section*{Management}
Treatment of acute adverse reactions is focused on symptomatic control and reassurance:
 
\begin{itemize}
    \item \textbf{Supportive environment:} placing the patient in a quiet, low-stimulus, dark room with gentle reassurance ("talking down").
    \item \textbf{Benzodiazepines:} first-line pharmacological treatment (e.g., diazepam) to control acute panic, severe anxiety, and mild hypertension/tachycardia.
    \item \textbf{Antipsychotics:} generally avoided unless severe agitation is refractory to benzodiazepines, as they can lower the seizure threshold or worsen anticholinergic symptoms.
    \item \textbf{HPPD management:} treated with lamotrigine, clonidine, or benzodiazepines, along with psychological support.
\end{itemize}
 
\section*{Complications}
\begin{itemize}
    \item Accidental injury or death due to impaired judgment and risky behavior
    \item Severe, persistent psychiatric distress (acute psychosis or exacerbation of underlying schizophrenia)
    \item Hallucinogen Persisting Perception Disorder (HPPD)
    \item Mild serotonin syndrome if combined with other serotonergic agents (like SSRIs or MAOIs)
\end{itemize}
 
\section*{Prognosis and Outcomes}
The prognosis for acute LSD intoxication is excellent. Direct physical toxicity is extremely low, and the drug is cleared from the body within 24 hours. Most patients recover fully without long-term physical sequelae once the drug's effects wear off. Psychological recovery depends on the presence of pre-existing psychiatric conditions; individuals with a history of psychosis may experience prolonged psychiatric symptoms. HPPD is rare, but can be chronic and disabling, requiring long-term specialist management.
 
\section*{Prevention and Self-Care}
\begin{itemize}
    \item The only reliable way to prevent adverse reactions or HPPD is to avoid the use of illicit LSD
    \item If someone chooses to use LSD, they should never do so alone; a sober, trusted "trip sitter" should be present
    \item Avoid combining LSD with alcohol, cannabis, stimulants, or prescription medications
    \item Ensure a safe, familiar, and calm environment before exposure to minimize the risk of a "bad trip"
    \item Seek immediate mental health support if experiencing persistent psychological distress, flashbacks, or mood changes post-use
\end{itemize}
 
\section*{Sources and Further Reading}
The following reputable organisations provide further information for healthcare students and patients seeking reliable, current advice about LSD.
 
\begin{itemize}
    \item Alcohol and Drug Foundation (ADF). \textit{Detailed information on LSD, its effects, risks, and support services.} https://adf.org.au/
    \item Healthdirect Australia. \textit{Hallucinogens: LSD effects and addiction support.} https://www.healthdirect.gov.au/
    \item Mayo Clinic. \textit{Drug addiction (substance use disorder): recovery and support.} https://www.mayoclinic.org/
    \item MSD Manual. \textit{Hallucinogen toxicity, symptoms, and clinical management.} https://www.msdmanuals.com/
    \item World Health Organization. \textit{Public health resources on psychoactive substances.} https://www.who.int/
\end{itemize}
